\documentclass[11pt, reqno]{amsart}

%%% Font Setup
% \usepackage{ETbb}
\usepackage{libertinus}           % text + math, type1 so pdflatex-safe
\usepackage[T1]{fontenc}

%%% Packages
\usepackage{amsmath, amssymb, amsthm}
\usepackage{xcolor}
\usepackage{bm}
\usepackage[colorlinks=true, linkcolor=blue!60!black, citecolor=blue!60!black, urlcolor=blue!60!black]{hyperref}
\usepackage{enumitem}
\usepackage{graphicx}
\usepackage{subcaption}
\usepackage{flafter}

\usepackage[framemethod=TikZ]{mdframed}
\mdfdefinestyle{lecturebox}{
  backgroundcolor=black!5,      % light gray fill
  linecolor=black!12,           % faint border; set linewidth=0pt to drop it
  linewidth=0.6pt,
  roundcorner=3pt,
  leftmargin=1em, rightmargin=1em,        % outer indent from text block
  innerleftmargin=1.4em, innerrightmargin=1.4em,
  innertopmargin=1em, innerbottommargin=1em,
}

%%% Page Geometry
\usepackage{geometry}
\geometry{margin=1.2in, top=1in, bottom=1in}

%%% Section Headings in Small Caps
\makeatletter
\def\section{\@startsection{section}{1}%
  \z@{.7\linespacing\@plus\linespacing}{.5\linespacing}%
  {\normalfont\scshape\centering}}
\makeatother

%%% Title in Small Caps
\makeatletter
\def\@settitle{\begin{center}%
  \baselineskip14\p@\relax
  \bfseries\scshape\uppercasenonmath\@title
  \@title
  \end{center}%
}
\makeatother

%%% Theorem Environments
\newtheoremstyle{notestyle}
  {6pt}{6pt}{\itshape}{}{\bfseries}{.}{.5em}
  {\thmname{#1}\thmnumber{ #2}\thmnote{ \normalfont(#3)}}

\theoremstyle{notestyle}
\newtheorem{theorem}{Theorem}[section]
\newtheorem{lemma}[theorem]{Lemma}
\newtheorem{proposition}[theorem]{Proposition}
\newtheorem{corollary}[theorem]{Corollary}

\newtheoremstyle{defstyle}
  {6pt}{6pt}{\normalfont}{}{\bfseries}{.}{.5em}
  {\thmname{#1}\thmnumber{ #2}\thmnote{ \normalfont(#3)}}

\theoremstyle{defstyle}
\newtheorem{definition}[theorem]{Definition}
\newtheorem{example}[theorem]{Example}
\newtheorem{remark}[theorem]{Remark}

\newcommand{\R}{\mathbb{R}}
\newcommand{\T}{^\top}
\newcommand{\mc}[1]{\mathcal{#1}}

%%% Additional notation
\newcommand{\E}{\mathbb{E}}
\newcommand{\V}{\mathcal{V}}
\newcommand{\Hphys}{\mathcal{H}_{\mathrm{phys}}}
\newcommand{\Cov}{\operatorname{Cov}}
\newcommand{\Tr}{\operatorname{Tr}}
\newcommand{\Id}{\mathrm{Id}}
\newcommand{\argmin}{\operatorname*{arg\,min}}
\newcommand{\norm}[1]{\left\lVert #1 \right\rVert}
\newcommand{\inner}[2]{\left\langle #1,#2 \right\rangle}
\newcommand{\MPS}{\mathcal{M}_{\chi}}
\newcommand{\Q}{\mathcal{Q}}
\newcommand{\FS}{\mathrm{FS}}
\newcommand{\Heff}{H_{\mathrm{eff}}}
\newcommand{\Neff}{N_{\mathrm{eff}}}

%%% Metadata
\title{Correlated Cubic Features}
\author{authors}
\date{}

\begin{document}
\maketitle
\thispagestyle{empty}

\noindent\textit{%
We need to extend Nisch's two-feature calculation to $k$ correlated target features at Bayesian equilibrium. In particular, when does one committee serve all the features, when does it split, and what populations of neurons minimize the posterior cost after splitting?}

% \tableofcontents

% \newpage

\section{Setup}

Let $N$ be the input dimension, $K$ be the number of hidden neurons, and $P$ be the number of training examples.

Take
\begin{equation}
    K = \frac{N^2}{\Delta},\quad P = \chi K,
\end{equation}
with $\Delta, \chi > 0$ fixed. Our sample ratio remains
\begin{equation}
    \alpha = \frac PN = \frac{\chi N}{\Delta},
\end{equation}
growing with $N$. The source's sample ratio is $\chi = P/K$, not our $\alpha$.

\section{Teacher and Training Data}
Draw independent $\mathbf x^\mu \sim \mathcal N(0, I_N)$. Choose unit teacher directions satisfying
\begin{equation}
    \mathbf e_\ell^\top \mathbf e_t = G_{\ell t} = (1 - \omega)d_{\ell t} + \omega,\quad 0 \leq \omega \leq 1.
\end{equation}
Initially use noiseless labels
\begin{equation}
    y^\mu = f_*(\mathbf x^\mu),\quad f_*(\mathbf x) = B\sum_{\ell = 1}^k \psi_3(\mathbf e_\ell^\top \mathbf x),\quad \psi_3(z) = \frac{H_3(z)}{\sqrt 6} = \frac{z^3 - 3z}{\sqrt{6}}.
\end{equation}
Where $H_3$ is the third Hermite polynomial. The normalization makes each feature have unit variance. For Gaussian inputs,
\begin{equation}
    \mathbb E_{\mathbf x} [\psi_3(\mathbf e_\ell^\top \mathbf x)\psi_3(\mathbf e_t^\top\mathbf x)] = G_{\ell t}^3.
\end{equation}
Thus $\omega$ is a direction overlap. The corresponding cubic-feature correlation is $\omega^3$.
\section{Student network and equilibrium ensemble}
Use the pure cubic student
\begin{equation}
    f_\theta(\mathbf x) = \sum_{i=1}^K a_i \gamma_3 \psi_3\left( \frac{\mathbf w_i^\top \mathbf x}{\sqrt N}\right),\quad \gamma_3 > 0.
\end{equation}
For this calculation, the prior is uniform over hidden-weight spheres and over readouts satisfying:
\begin{equation}
    \|\mathbf w_i\|^2 = N,\quad |a_i| \leq \frac{A}{\sqrt K},\quad \sum_{i=1}^K a_i^2 = \tau^2,\quad 0 < \tau^2 < A^2.
\end{equation}
Here, $A, \tau, \gamma_3$ are fixed. The factor $K^{-1/2}$ in Nisch's network has been absorbed into our readouts $a_i$. For fixed dataset $\mathcal D$, define
\begin{equation}
    E_{\mathcal D}(\theta) = \frac 12 \sum_{\mu=1}^P[y^\mu - f_\theta(\mathbf x^\mu)]^2,\quad p_\beta(\theta \mid \mathcal D) = \frac{p_0(\theta)e^{-\beta E_{\mathcal D}(\theta)}}{Z_\beta (\mathcal D)}.
\end{equation}
This is an equilibrium problem. It doesn't specify how quickly SGD or a posterior sampler reaches the predicted populations. We write $\langle \cdot \rangle_{\mathcal D}$ for posterior averages and $\mathbb E_{\mathbf x}$ for averages over fresh Gaussian inputs.

Define $\mathbf w_\ell^* = \sqrt N \mathbf e_\ell$. The signed teacher overlaps are
\begin{equation}
    m_{i\ell} = \frac{\mathbf w_i^\top \mathbf w_\ell^*}{N},\quad q_i = \frac{\|\mathbf w_i\|^2}{N} = 1.
\end{equation}
\section{Reduced description in terms of committees}
Let $\mathcal V = \mathrm{span}\{\mathbf e_1, \ldots, \mathbf e_k\}$. A committee consists of neurons sharing a projection into $\mathcal V$:
\begin{equation}
    \frac{\mathbf w_i}{\sqrt N} = r_c \mathbf u_c + \sqrt{1 - r_c^2 }\bm \xi_i,\quad \mathbf u_c \in \mathcal V,\quad \|\mathbf u_c\| = \|\bm \xi_i\| = 1,\quad \bm\xi_i \perp \mathcal V.
\end{equation}
$\mathcal V$ is a small subspace inside the input space such that the teacher depends only on the projections of $\mathbf x$ into $\mathcal V$. We can therefore split a neuron's normalized weight into the above decomposition. The first term is the projection into the teacher subspace, the second term is the transverse component. Here, $\mathbf u_c$ tells us which direction the group represents, while $r_c$ tells us how \textit{strongly neurons concentrates on that directions.} A committee shares the first component; its neurons can differ in their transverse components. The teacher overlaps are $m_{i\ell} = r_c \mathbf u_c^\top \mathbf e_\ell$.

The useful identity here is
\begin{equation}
    \mathbb E_Z \left[ \psi_3 (rz + \sqrt{1-r^2}Z)\right] = r^3 \psi_3(z).
\end{equation}
So part of the neuron's response that survives averaging over transverse inputs is a cubic feature along $\mathbf u_c$, with strength proportional to $r_c^3$. 

Under the spherical prior, most weight distributions are almost orthogonal to any fixed low-dim subspace. A finite projection $r$ is rare because comparitively few directions have that projection:
\begin{equation}
    \mathrm{Pr}(r) \asymp e^{-N J(r)},\quad J(r) = -\frac 12\log(1 - r^2).
\end{equation}
This is an \textit{entropy cost}. Every allowed direction is equally likely, but there are vastly more directions with small project than large.

Stronger alignment gives more cubic output per neuron (prop to $r^3$), but costs more prior probability. 

Nisch' leading-order calculation minimizes the alignment cost per unit cubic output:
\begin{equation}
    \frac{\text{alignment cost}}{\text{cubic output}} \propto \frac{J(r)}{r^3},\quad r_* = \argmin_{0 < r < 1} \frac{J(r)}{r^3} \simeq 0.73047.
\end{equation}
So weak alignment produces too little signal; alignment approaching one becomes too statistically costly. 

After optimizing readout magnitudes and projection lengths, represent the retained cubic signal as
\begin{equation}
    f_{\mathcal C}(\mathbf x) = \sum_{c=1}^L b_c \psi_3(\mathbf u_c^\top \mathbf x),\quad b_c \geq 0.
\end{equation}
\textit{The number $L$ of committees is unknown.} Signs can also be absorbed into $\mathbf u_c$, since $\psi_3$ is odd.

Under the background approximation, the effective population cost is
\begin{equation}
    F(\mathcal C) = \lambda \sum_c b_c + \frac 12 \mathbb E_{\mathbf x}\left[ f_*(\mathbf x) - f_{\mathcal C}(\mathbf x)\right]^2
\end{equation}
with 
\begin{equation}
    \lambda = \frac{\sqrt{\Delta}J(r_*)}{g_{\text{eff}}\gamma_3 A r_*^3}.
\end{equation}

\section{Three Correlated Features}

Similarly to the two correlated features notes, let
\begin{align*}
f_* (x) = \beta \sum_{a=1}^3 \psi_3(e_a^\top x),
\qquad e_a^\top e_b = \omega \quad (a \neq b)
\end{align*}
Thus, we have three features, with mutual inner products all of $\omega$.
This case will be easier than the four-feature asymmetric case; not only do we have one fewer feature, there is also an $S_3$ symmetry.

We can first rederive the free energy expression here, in a very similar way to the two feature case.
If $w$ is one hidden weight, we denote by $p$ its projection into $S = e_1 \times e_2 \times e_3$ space, and $u$ the unit direction of that projection:
\begin{align*}
w = pu + \sqrt{1-p^2} w_\perp
\end{align*}
As in the two-feature calculation,
\begin{align*}
\mathbb{E}_{x_\perp} \psi_3(w^\top x)
= p^3 \psi_3 (u^\top x).
\end{align*}
We again have that we need to minimize the alignment cost, which requires minimizing
\begin{align*}
\frac{J(p)}{p},
\qquad J(p) = -\frac{1}{2} \log (1-p^2)
\end{align*}
This all results in the same preferred radius $p = r_3 \approx 0.73$, and the same optimal $\lambda$.
Thus, these values do not depend on the number of features.

The ``free energy" that ends up being minimized is written
\begin{align*}
F = \lambda \sum_{j} b_j + \frac{1}{2} \mathbb{E}_x [f_*(x) - f_C(x)]^2
\qquad f_C(x) = \sum_j b_j \psi_3(u_j^\top x)
\end{align*}
as before.

Now, we perform a similar centering step as in the two-feature case, this time around a central axis, in light of the $S_3$ symmetry.
Define the normalized symmetric axis vector
\begin{align*}
e_0 = \frac{e_1 + e_2 + e_3}{\sqrt{3(1+2\omega)}}
\end{align*}
which lets us write any of the teacher vectors as
\begin{align*}
    e_a = c e_0 + s v_a,
    \qquad c = \sqrt{\frac{1+2\omega}{3}},
    \qquad s = \sqrt{\frac{2(1-\omega)}{3}}
\end{align*}
where the $v_a$ are three unit vectors in the plane perpendicular to $e_0$, each pointing towards $e_a$, which form an equilateral triangle in that plane.

We can parametrize the projection of $w$, a committee representative, as
\begin{align*}
    w = pu + \sqrt{1-p^2}w_\perp,
    \qquad u = t e_0 + \sqrt{1-t^2} v
\end{align*}
with $v$ lying in the $e_0$ perpendicular plane.

Now, similarly to the two feature case, we look at the difference in free energy from adding a single committee with output $b \psi_3 (u^\top x)$ to a completely unaligned network with $f_C = 0$:
\begin{align*}
    F(b, u) - F(0)
    &= b \left[\lambda - \mathbb{E}_x \{f_* \psi_3 (u^\top x)\} \right]
    + \frac{b^2}{2} \mathbb{E}_x \psi_3 (u^\top x)^2 \\
    &= b \left[ \lambda - \beta \{ (e_1^\top x)^3 + (e_2^\top x)^3 + (e_3^\top x)^3 \} 
    \right] + \frac{b^2}{2}
\end{align*}
where we used Hermite orthogonality to evaluate the expectations.
Then define
\begin{align*}
    H(t, \theta) = \sum_{a=1}^3 (e_a^\top u)^3
\end{align*}
where we have written on the left side the dependence of $H$ on $t$ and $\theta$, which we define as the angle between $v$ and $v_1$.
After some algebraic magic happens, it can be shown that
\begin{align}\label{eq:three-feature-H}
    H(t, \theta)
    &= 3c^3 t^3
    + \frac{9}{2} cs^2 t (1-t^2)
    + \frac{3}{4} s^3 (1-t^2)^{3/2} \cos (3\theta)
\end{align}
Recall that $e_a = c e_0 + s v_a$ and $u = t e_0 + \sqrt{1-t^2}v$.
Schematically, then, each of these terms in $H$ comes from the different combinations of terms in $e_a$ and $u$.
The first corresponds to three $e_0$ components, the second to one $e_0$ and two transverse components.
It turns out the term from two $e_0$ and one transverse vanishes because $v_1 + v_2 + v_3 = 0$. 
The last term is the pure transverse term.
The $\cos 3\theta$ is the lowest angular harmonic that is compatible with the $S_3$ symmetry of the teacher vectors: we must have $H(t, \theta + 2\pi/3) = H(t, \theta)$.

So, we can write the change in free energy as
\begin{align*}
    \Delta F
    = b \left[\lambda - \beta H(u) \right] + \frac{b^2}{2}
\end{align*}
For $0 < b \ll 1$, adding the committee will be favored when maximizing $H(u)$, as in the two-feature case.
Examining the formula \ref{eq:three-feature-H}, we see that $H$ is maximized for $\cos 3\theta = 1$, which gives three equivalent directions, corresponding to $e_1$, $e_2$, and $e_3$.

Then we can check when the various solutions are favored.
$t = 1$ corresponds to the direction of the committee being completely shared, giving $H_\text{sh} = 3c^3$.
We can check when this solution is locally stable.
Let $q = \sqrt{1-t^2} \ll 1$.
Then,
\begin{align*}
    H
    &= 3c^3 (1-q^2)^{3/2}
    + \frac{9}{2} cs^2 (1-q^2)q^2
    + \frac{3}{4} s^3 q^3 \cos (3\theta) \\
    &= 3c^3 + \frac{9}{2} c(s^2 - c^2)q^2
    + \frac{3}{4} s^3 q^3 \cos 3\theta + O(q^4)
\end{align*}
The $q^2$ term is negative (for recall we're trying to \emph{maximize} $H$), and this solution is stable, when $s^2 < c^2$, which is exactly when $\omega > 1/4$.

However, it turns out this is not necessarily the globally-stable point for $\omega > 1/4$.
We also need to consider the asymmetric solutions.
WLOG, let
\begin{align*}
(e_1^\top u, e_2^\top u, e_3^\top u)
= (x, y, y), \qquad x > y
\end{align*}
Then we want to maximize $H(u) = x^3 + 2y^3$.
To globally solve this constrained maximization problem, it helps to change coordinate systems.
Let
\begin{align}
    u = \sum_b \alpha_b e_b
\end{align}
If $G_{ab} = e_a^T e_b$ is the Gram matrix, we can write the projection of $u$ onto each feature vector as
\begin{align*}
    r_a  = e_a^\top u
    = \sum_b G_{ab} \alpha_b,
    \qquad \text{or }
    r &= G \alpha \\
    \implies
    \alpha &= G^{-1} r
\end{align*}
Thus, we can write the constraint that $u$ is a unit vector as
\begin{align*}
    |u|^2 = \alpha^\top G \alpha = r^\top G^{-1}  r = 1
\end{align*}
In this case with equally-correlated features,
\begin{align*}
    G
    = \begin{pmatrix}
        1 & \omega & \omega \\
        \omega & 1 & \omega \\
        \omega & \omega & 1
    \end{pmatrix}
    &= (1-\omega)I + \omega 1 1^\top \\
    G^{-1}
    &= \frac{1}{1-\omega} \left[I - \frac{\omega}{1 + 2\omega} 1 1^\top \right]
\end{align*}
We can now write the constrained maximization problem using a Lagrange mulitiplier $\mu$ as
\begin{align*}
    \mathcal{L} = \sum_a r_a^3 - \mu (r^\top G^{-1} r - 1)
\end{align*}
At the max, we have $\partial \mathcal{L} / \partial r_a = 0$, which gives
\begin{align}\label{eq:H-stationarity-lagrange}
    3r_a^2 = 2\mu (G^{-1} r)_a
\end{align}
We write the expression on the RHS as
\begin{align*}
    (G^{-1} r)_a
    = \frac{1}{1-\omega} \left[ r_a - \frac{\omega}{1 + 2\omega} (x + 2y) \right]
\end{align*}
Explicitly, the first and the symmetric components are
\begin{align*}
    (G^{-1} r)_1
    &= \frac{1}{1-\omega} \left[ x - \frac{\omega}{1 + 2\omega} (x + 2y) \right] \\
    &= \frac{(1 + \omega)x - 2\omega y}{(1-\omega)(1 + 2\omega)} \\
    (G^{-1}r)_{2,3}
    &= \frac{1}{1-\omega} \left[ y - \frac{\omega}{1 + 2\omega} (x + 2y) \right] \\
    &= \frac{y - \omega x}{(1-\omega)(1+2\omega)}
\end{align*}
Plugging these into equation \ref{eq:H-stationarity-lagrange}, we get the stationarity equations
\begin{align*}
    3x^2
    &= 2\mu \frac{(1 + \omega)x - 2\omega y}{(1-\omega)(1 + 2\omega)} \\
    3y^2
    &= 2 \mu \frac{y - \omega x}{(1-\omega)(1+2\omega)}
\end{align*}
Dividing and letting $R = x/y$, we have
\begin{align*}
    R^2
    &= \frac{(1+\omega)R - 2\omega}{1 - \omega R} \\
    \implies
    0 &= (R - 1) \left[\omega R^2 - (1 - \omega)R + 2\omega \right]
\end{align*}
$R=1$ is the $x = y$ symmetric case, but we're interested in the asymmetric one, so we can divide that out, giving the final stationarity equation
\begin{align}\label{eq:stationarity}
    0 &= \omega R^2 - (1-\omega)R + 2\omega \\
    \implies
    \omega &= \frac{R}{R^2 + R + 2}
\end{align}
This equation will be satisfied exactly when the global maximum of $H$ is asymmetric.

Recall that for the symmetric case, $H_\text{sym}^2 = 3^2c^6 = 1/3 \times (1+2\omega)^3$, so
\begin{align*}
    \left( \frac{H_\text{asym}}{H_\text{sym}} \right)^2
    &= \frac{3y^6(R^3+2)^2}{(1+2\omega)^3}
\end{align*}
Then we need to replace $y$ and $\omega$ with $R$.
We can write $r = y(R, 1, 1)$, which lets us rewrite the $|u|^2 = 1$ constraint as
\begin{align*}
    1 &= \frac{y^2}{1-\omega} \left[ R^2 + 2 - \frac{\omega}{1+2\omega} (R+2)^2 \right] \\
    \implies y^2
    &= \frac{(1-\omega)(1+2\omega)}{(1+\omega)R^2 - 4\omega R +2}
\end{align*}
We'll let the denominator $D = (1+\omega)R^2 - 4\omega R + 2$, so
\begin{align*}
    y^6
    &= \frac{(1-\omega)^3 (1+2\omega))^3}{D^3} \\
    \implies
    \left( \frac{H_\text{asym}}{H_\text{sh}} \right)^2
    &= \frac{3(1-\omega)^3 (R^3+2)^2}{D^3}
\end{align*}
Then using equation \ref{eq:stationarity},
\begin{align*}
    1-\omega
    &= \frac{R^2+2}{R^2+R+2} \\
    D
    &= \frac{(R+2)(R^3+3)}{R^2+R+2}
\end{align*}
Plugging these into the above, we get
\begin{align*}
    \left( \frac{H_\text{asym}}{H_\text{sh}} \right)^2
    &= \frac{3(R^2+2)^3}{(R+2)^3(R^3+2)}
\end{align*}
The crossover point in favorability here is when the numerator and denominator are equal, which we can find:
\begin{align*}
    0
    &= 3(R^2+2)^2 - (R+2)^3(R^3+2)
    = 2(R-1)^3 (R^3-4)
\end{align*}
So, they are equally favored when (trivially) $R=1$, and nontrivially when $R = 2^{2/3}$.
Plugging this into equation \ref{eq:stationarity} gives
\begin{align*}
    \omega_c
    &= \frac{2^{2/3}}{2^{4/3} + 2^{2/3} + 2}
    = 2^{1/3} - 1
    \approx 0.259921 > \frac{1}{4}
\end{align*}
How large can $\omega$ get and keep the asymmetric solution locally stable?
Maximize equation \ref{eq:stationarity}:
\begin{align*}
    \frac{d\omega}{dR}
    &= \frac{(R^2 + R + 2) - R(2R+2)}{(R^2+R+2)^2} \\
    &= \frac{2-R^2}{(R^2+R+2)^2}
\end{align*}
So the max is at $R = 2^{1/2}$, and
\begin{align*}
    \omega_\text{spin}
    &= \frac{\sqrt{2}}{4+\sqrt{2}}
    \approx 0.261204
\end{align*}
So for correlations above $\omega_\text{spin}$, the shared feature is global free energy minimum, and the only local minimum. At $\omega_c = 2^{1/3}-1$, the asymmetric solution becomes a local minimum, but the shared feature is still global minimum.
Then between $\omega_c$ and $\omega = 1/4$, the symmetric solution remains locally-stable, but the asymmetric solution is the global minimum.
And below $\omega = 1/4$, the symmetric solution loses local stability, and the asymmetric solution is the only stable one.

\section{Spherical Harmonics for Any Number of Cubic Features}
\label{sec:cubic-spherical-harmonics}

The angular structure of $H$ in the three-feature calculation has a general explanation.
For $k$ unit teacher directions, the same Hermite overlap identity gives
\begin{align*}
    f_*(\mathbf x) &= B\sum_{\ell=1}^k \psi_3(\mathbf e_\ell^\top\mathbf x), \\
    \mathbb E_{\mathbf x}\left[f_*(\mathbf x)\psi_3(\mathbf u^\top\mathbf x)\right]
    &= B H(\mathbf u),
    \qquad H(\mathbf u) = \sum_{\ell=1}^k (\mathbf e_\ell^\top\mathbf u)^3.
\end{align*}
Here, $\mathbf u$ is a unit direction in the teacher span $\mathcal V$.
Thus, adding a committee to the unaligned state still changes the cost by
\begin{align*}
    F(b,\mathbf u)-F(0) = b[\lambda-BH(\mathbf u)]+\frac{b^2}{2}.
\end{align*}
We now allow arbitrary correlations between the teacher directions, rather than requiring all pairwise overlaps to equal $\omega$.
Let $m=\dim\mathcal V$ and use orthonormal coordinates in this span.
The angular variable then lies on $S^{m-1}$; for linearly independent teachers, $m=k$.
This distinction matters because the raw teacher overlaps describe an ellipsoid when the teachers are correlated, not a Euclidean unit sphere.

A generalized spherical harmonic of degree $j$ is the restriction to the unit sphere of a homogeneous degree-$j$ polynomial whose Euclidean Laplacian vanishes.
We will show that $H$ contains only degrees $1$ and $3$.
Write these two parts as $H^{(1)}$ and $H^{(3)}$, to distinguish them from the Hermite polynomial $H_3$ in the definition of $\psi_3$.

To find the decomposition, first extend $H$ off the sphere:
\begin{align*}
    P(\mathbf z) = \sum_{\ell=1}^k (\mathbf e_\ell^\top\mathbf z)^3,
    \qquad \mathbf v = \sum_{\ell=1}^k \mathbf e_\ell,
    \qquad \mathbf z\in\mathcal V.
\end{align*}
This polynomial is homogeneous of degree three, but it is generally not harmonic.
Since every teacher direction has unit norm,
\begin{align*}
    \Delta_{\mathbf z}(\mathbf e_\ell^\top\mathbf z)^3
    &= 6\|\mathbf e_\ell\|^2(\mathbf e_\ell^\top\mathbf z)
    = 6\mathbf e_\ell^\top\mathbf z, \\
    \Delta_{\mathbf z}P(\mathbf z) &= 6\mathbf v^\top\mathbf z.
\end{align*}
Here $\Delta_{\mathbf z}$ denotes the Laplacian in the $m$ orthonormal coordinates of $\mathcal V$, not the sample-scaling parameter $\Delta$.
The term we need to subtract follows from the product rule:
\begin{align*}
    \Delta_{\mathbf z}\left(\|\mathbf z\|^2\mathbf v^\top\mathbf z\right)
    &= 2m\mathbf v^\top\mathbf z+4\mathbf v^\top\mathbf z \\
    &= 2(m+2)\mathbf v^\top\mathbf z.
\end{align*}
Therefore, the polynomials
\begin{align*}
    Q_1(\mathbf z) &= \frac{3}{m+2}\mathbf v^\top\mathbf z, \\
    Q_3(\mathbf z) &= P(\mathbf z)-\|\mathbf z\|^2Q_1(\mathbf z)
\end{align*}
are homogeneous of degrees one and three, respectively, and both have zero Laplacian.
Restricting to $\|\mathbf u\|=1$ gives the desired decomposition:
\begin{align}
    H(\mathbf u) &= H^{(1)}(\mathbf u)+H^{(3)}(\mathbf u),
    \label{eq:H-harmonic-decomposition} \\
    H^{(1)}(\mathbf u) &= \frac{3}{m+2}\mathbf v^\top\mathbf u,
    \label{eq:H-degree-one} \\
    H^{(3)}(\mathbf u)
    &= \sum_{\ell=1}^k\left[(\mathbf e_\ell^\top\mathbf u)^3
    -\frac{3}{m+2}\mathbf e_\ell^\top\mathbf u\right].
    \label{eq:H-degree-three}
\end{align}
For $m\geq2$, these are spherical harmonics of degrees one and three on $S^{m-1}$.
Equivalently, expanding $H$ in any complete spherical-harmonic basis produces coefficients only at these two degrees.
The decomposition is exact, regardless of the correlations or the number of teachers.
If the span is one-dimensional, the sphere consists of just two points and $P(\mathbf z)=\|\mathbf z\|^2Q_1(\mathbf z)$, so $Q_3$ vanishes.

We can also express the answer using harmonics centered on each teacher direction.
For $m\geq2$, define
\begin{align*}
    Z_1^{(m)}(t)=t,
    \qquad Z_3^{(m)}(t)=\frac{(m+2)t^3-3t}{m-1}.
\end{align*}
The functions $Z_j^{(m)}(\mathbf e_\ell^\top\mathbf u)$ are zonal spherical harmonics: they depend only on the angle from $\mathbf e_\ell$.
Their normalization is $Z_j^{(m)}(1)=1$.
For $m>2$, they are the normalized Gegenbauer polynomials with parameter $(m-2)/2$; for $m=2$, they satisfy $Z_j^{(2)}(\cos\theta)=\cos(j\theta)$ for $j=1,3$.
Rearranging their definitions gives
\begin{align}
    H(\mathbf u)
    = \frac{3}{m+2}\sum_{\ell=1}^k Z_1^{(m)}(\mathbf e_\ell^\top\mathbf u)
    +\frac{m-1}{m+2}\sum_{\ell=1}^k Z_3^{(m)}(\mathbf e_\ell^\top\mathbf u).
    \label{eq:H-zonal-harmonics}
\end{align}

For the three-feature geometry above, $m=3$ when $\omega<1$, and $\mathbf v=3ce_0$.
The degree-one part is therefore $H^{(1)}(t,\theta)=9ct/5$.
Using $c^2+s^2=1$, equation \eqref{eq:three-feature-H} can be rewritten as
\begin{align*}
    H(t,\theta)
    &= \frac{9c}{5}t
    +\frac{3c}{5}(5c^2-3)P_3(t)
    +\frac{3}{4}s^3(1-t^2)^{3/2}\cos(3\theta), \\
    P_3(t)&=\frac{5t^3-3t}{2}.
\end{align*}
Here $P_3$ is the third Legendre polynomial.
The first term is a degree-one spherical harmonic, while the last two are degree-three spherical harmonics.
Thus, the $\cos(3\theta)$ term found earlier is one part of the full degree-three angular dependence.

The same proof also allows unequal feature strengths: replacing each summand in $H$ by $\eta_\ell(\mathbf e_\ell^\top\mathbf u)^3$ replaces $\mathbf v$ by $\sum_\ell\eta_\ell\mathbf e_\ell$ and weights each summand in \eqref{eq:H-degree-three} by $\eta_\ell$.
The degree-one component vanishes exactly when this weighted sum of directions is zero.
Although the target consists of cubic Hermite features in the Gaussian input, its overlap with a candidate committee generally contains both degree-one and degree-three spherical harmonics as a function of the committee direction.

\appendix
\section{Generalized Spherical Harmonics}
\label{app:generalized-spherical-harmonics}

Ordinary spherical harmonics describe angular dependence on $S^2$, the unit sphere in $\R^3$.
Here, ``generalized'' means that we allow the ambient space to be $\R^m$ and the sphere to be
\begin{align*}
    S^{m-1}=\{\mathbf u\in\R^m:\|\mathbf u\|=1\}.
\end{align*}
These functions are also called hyperspherical harmonics.
They play the same role as the familiar $Y_{\ell q}(\theta,\phi)$: they separate angular dependence into orthogonal modes of increasing degree.
We use $m$ for the dimension and $q$ for a mode label, so $m$ is not the magnetic quantum number of the usual three-dimensional notation.
In our problem, $m=\dim\mathcal V$, the dimension of the teacher span, rather than the input dimension $N$.
We assume $m\geq2$ throughout this appendix.

\subsection{From harmonic polynomials to angular modes}
One way to recognize an ordinary spherical harmonic is to multiply it by the appropriate radial power: $r^\ell Y_{\ell q}$ is a homogeneous harmonic polynomial in three Cartesian coordinates.
The same construction works in any dimension.
Let $Q_\ell$ be a polynomial satisfying
\begin{align*}
    Q_\ell(r\mathbf u)=r^\ell Q_\ell(\mathbf u),
    \qquad \Delta_{\mathbf z}Q_\ell(\mathbf z)=0.
\end{align*}
The first condition says that every monomial has total degree $\ell$; the second says that the polynomial is harmonic.
Its restriction $Y(\mathbf u)=Q_\ell(\mathbf u)$ to $S^{m-1}$ is a generalized spherical harmonic of degree $\ell$.
The polynomial $Q_\ell$ contains both radial and angular dependence, while $Y$ contains only the angular part.

To connect this definition to the familiar angular eigenvalue equation, write $\mathbf z=r\mathbf u$.
The Euclidean Laplacian separates as
\begin{align*}
    \Delta_{\mathbf z}
    =\frac{\partial^2}{\partial r^2}
    +\frac{m-1}{r}\frac{\partial}{\partial r}
    +\frac{1}{r^2}\Delta_{S^{m-1}},
\end{align*}
where $\Delta_{S^{m-1}}$ is the angular Laplacian, or Laplace--Beltrami operator, on the unit sphere.
Applying it to $Q_\ell(r\mathbf u)=r^\ell Y(\mathbf u)$ gives
\begin{align*}
    0=\Delta_{\mathbf z}Q_\ell
    =r^{\ell-2}\left[\ell(\ell+m-2)Y+\Delta_{S^{m-1}}Y\right].
\end{align*}
Thus,
\begin{align}
    -\Delta_{S^{m-1}}Y=\ell(\ell+m-2)Y.
    \label{eq:appendix-harmonic-eigenvalue}
\end{align}
For $m=3$, this is exactly the usual eigenvalue $\ell(\ell+1)$.
For $m=2$, the sphere is a circle, the angular Laplacian is $\partial_\theta^2$, and the positive-degree modes are $\cos(\ell\theta)$ and $\sin(\ell\theta)$.

\subsection{Orthogonality and the number of modes}
As in three dimensions, the degree does not specify a unique harmonic.
The dimension of the degree-$\ell$ harmonic space is
\begin{align*}
    d_{m,\ell}
    =\binom{m+\ell-1}{\ell}-\binom{m+\ell-3}{\ell-2},
\end{align*}
where the second term is taken to be zero for $\ell<2$.
This gives $d_{3,\ell}=2\ell+1$, the usual number of magnetic modes.
On the circle, there is one constant mode and two modes at each positive degree.
We can choose a real orthonormal basis $Y_{\ell q}$, with $q=1,\ldots,d_{m,\ell}$, such that
\begin{align*}
    \int_{S^{m-1}}Y_{\ell q}(\mathbf u)Y_{\ell' q'}(\mathbf u)\,d\Omega(\mathbf u)
    =\delta_{\ell\ell'}\delta_{qq'}.
\end{align*}
Here $d\Omega$ is the usual surface-area measure.
Different degrees are orthogonal because the angular Laplacian is self-adjoint and their eigenvalues in \eqref{eq:appendix-harmonic-eigenvalue} are distinct.
Together, these modes form a complete basis for square-integrable functions on the sphere, just as ordinary spherical harmonics do on $S^2$.

\subsection{Zonal harmonics and Gegenbauer polynomials}
For our overlap function, a useful harmonic is one centered on a unit direction $\mathbf e$ and depending only on $t=\mathbf e^\top\mathbf u$, the cosine of the angle from that direction.
Such a function is called a zonal harmonic; on $S^2$, the familiar example is $P_\ell(\cos\theta)$ with the polar axis along $\mathbf e$.
For a function $Z(t)$, the angular eigenvalue equation becomes
\begin{align*}
    (1-t^2)Z''(t)-(m-1)tZ'(t)+\ell(\ell+m-2)Z(t)=0.
\end{align*}
For $m>2$, its degree-$\ell$ polynomial solution, normalized at $t=1$, is
\begin{align*}
    Z_\ell^{(m)}(t)
    =\frac{C_\ell^{(\nu)}(t)}{C_\ell^{(\nu)}(1)},
    \qquad \nu=\frac{m-2}{2},
\end{align*}
where $C_\ell^{(\nu)}$ is a Gegenbauer polynomial.
Thus Gegenbauer polynomials replace Legendre polynomials when the dimension changes.
At $m=3$, $\nu=1/2$ and $Z_\ell^{(3)}=P_\ell$.
At $m=2$, we instead use $Z_\ell^{(2)}(\cos\theta)=\cos(\ell\theta)$, avoiding the degenerate Gegenbauer normalization at $\nu=0$.
In particular,
\begin{align*}
    Z_1^{(m)}(t)=t,
    \qquad Z_3^{(m)}(t)=\frac{(m+2)t^3-3t}{m-1},
\end{align*}
which are the functions used in equation \eqref{eq:H-zonal-harmonics}.
A zonal harmonic about a chosen axis is a particular linear combination of the degree-$\ell$ basis functions, rather than the entire degree-$\ell$ space.

\subsection{Why a cubic has two harmonic degrees}
Polynomial degree and angular harmonic degree need not coincide.
For example, if $Q_1$ is linear, then $\|\mathbf z\|^2Q_1(\mathbf z)$ is a homogeneous cubic, but on the unit sphere it reduces to the degree-one harmonic $Q_1(\mathbf u)$.
More generally, a homogeneous polynomial of degree $p$ has a unique decomposition
\begin{align*}
    P_p(\mathbf z)
    =\sum_{j=0}^{\lfloor p/2\rfloor}\|\mathbf z\|^{2j}Q_{p-2j}(\mathbf z),
\end{align*}
where each $Q_{p-2j}$ is homogeneous and harmonic of the indicated degree.
Restricting to the sphere removes the radial factors, leaving angular degrees $p,p-2,\ldots$.
For $p=3$, this is precisely $P_3=Q_3+\|\mathbf z\|^2Q_1$, whose explicit construction appears in Section~\ref{sec:cubic-spherical-harmonics}.
This explains why cubic teacher features lead to degrees one and three in $H$, even when there are many features or their directions are correlated.

For further background, see Frye and Efthimiou, \href{https://arxiv.org/abs/1205.3548}{\textit{Spherical Harmonics in $p$ Dimensions}}, especially Chapter~4, and the \href{https://dlmf.nist.gov/18.38}{NIST discussion of zonal spherical harmonics}.

\end{document}
